\documentclass[../DoAn.tex]{subfiles}

\begin{document}

\section{Tổng quan chương}

Chương này khảo sát bài toán truy xuất và hỏi đáp thông tin trên tài liệu PDF,
đồng thời xác định các yêu cầu cần thiết cho hệ thống được xây dựng trong đồ án.
Khác với Chương 1, phần này không nhắc lại toàn bộ động cơ của đề tài mà tập
trung làm rõ đầu vào, đầu ra, nhóm người dùng, các thách thức chính và tiêu chí
đánh giá. Kết quả của chương là danh sách yêu cầu chức năng, yêu cầu phi chức
năng và bảng liên kết giữa yêu cầu với các phần thiết kế, thực nghiệm ở các
chương sau.

\section{Bài toán truy xuất và hỏi đáp trên tài liệu PDF}
\label{section:pdf-qa-problem}

Bài toán được xem xét trong đồ án có đầu vào là một hoặc nhiều tài liệu PDF và
câu hỏi của người dùng bằng ngôn ngữ tự nhiên. Đầu ra mong muốn là câu trả lời
ngắn gọn, đúng trọng tâm, kèm dẫn chứng cho phép người dùng kiểm tra lại nội
dung trong tài liệu gốc. Dẫn chứng có thể trỏ tới trang, đoạn văn, bảng, hàng hoặc
ô bảng tùy theo loại thông tin được sử dụng để trả lời.

Ở mức tổng quát, bài toán gồm hai nhóm nhiệm vụ. Nhóm thứ nhất là xử lý và tổ
chức tài liệu PDF thành các đơn vị có thể truy xuất, bao gồm văn bản, bảng và các
thông tin nguồn đi kèm. Nhóm thứ hai là truy xuất các đoạn thông tin liên quan,
kiểm tra mức độ đầy đủ của bằng chứng và tạo câu trả lời dựa trên các bằng chứng
đó. Nếu bước truy xuất lấy sai hoặc thiếu bằng chứng, bước hỏi đáp phía sau rất
khó tạo ra câu trả lời chính xác. Vì vậy, đồ án không chỉ đánh giá câu trả lời cuối
cùng mà còn đánh giá từng tầng của quy trình xử lý.

\begin{table}[H]
\centering
\small
\renewcommand{\arraystretch}{1.2}
\caption{Mô tả đầu vào, đầu ra và yêu cầu chính của bài toán}
\label{tab:problem-input-output}
\begin{tabularx}{\textwidth}{
    |>{\raggedright\arraybackslash}p{0.24\textwidth}
    |>{\raggedright\arraybackslash}X|
}
\hline
\textbf{Thành phần} & \textbf{Mô tả} \\
\hline
Đầu vào tài liệu &
Tài liệu PDF có lớp văn bản, tài liệu có bảng, bố cục nhiều vùng hoặc một số
trang cần nhận dạng ký tự từ ảnh. \\
\hline
Đầu vào truy vấn &
Câu hỏi bằng ngôn ngữ tự nhiên, có thể hỏi về đoạn văn, điều khoản, quy định,
bảng hoặc giá trị trong ô bảng. \\
\hline
Đầu ra truy xuất &
Danh sách bằng chứng liên quan kèm thông tin nguồn như tài liệu, trang, loại
nội dung và vị trí trong tài liệu. \\
\hline
Đầu ra hỏi đáp &
Câu trả lời ngắn gọn, bám sát bằng chứng và có dẫn chứng để kiểm chứng. \\
\hline
Yêu cầu kiểm soát &
Khi tài liệu không chứa đủ thông tin, hệ thống cần trả lời thận trọng thay vì tự
suy đoán. \\
\hline
\end{tabularx}
\end{table}

\section{Đặc điểm dữ liệu PDF và các thách thức chính}
\label{section:pdf-data-characteristics}

PDF là định dạng thuận tiện cho hiển thị, lưu trữ và chia sẻ, nhưng không phải lúc
nào cũng lưu nội dung theo cấu trúc thuận lợi cho xử lý tự động. Nhiều tài liệu PDF
chứa văn bản, bảng, hình ảnh, chú thích, tiêu đề đầu trang, chân trang, số trang và
bố cục nhiều cột. Khi trích xuất tự động, nội dung có thể bị sai thứ tự đọc, mất ngữ
cảnh hoặc lẫn các thành phần không thuộc nội dung chính. Các bộ dữ liệu phân
tích bố cục tài liệu như DocLayNet và PubTables-1M cũng cho thấy xử lý layout và
bảng là những bài toán riêng trong khai thác tài liệu PDF
\cite{pfitzmann2022doclaynet,smock2022pubtables}.

Bảng biểu là một thách thức quan trọng vì thông tin trong bảng phụ thuộc vào quan
hệ giữa hàng, cột và ô. Nếu bảng chỉ được nối thành văn bản tuyến tính, hệ thống có
thể giữ được một phần từ khóa nhưng mất quan hệ giữa điều kiện và giá trị. Điều
này ảnh hưởng trực tiếp đến các câu hỏi dạng tra cứu như quy đổi điểm, tìm khoảng
giá trị hoặc xác định điều kiện áp dụng. Với PDF được quét từ ảnh, hệ thống còn
cần nhận dạng ký tự quang học để lấy nội dung và vị trí từ trên trang
\cite{du2020ppocr}.

Từ các đặc điểm trên, đồ án tập trung vào bốn thách thức chính: trích xuất nội dung
PDF theo đúng thứ tự đọc, biểu diễn bảng sao cho giữ được cấu trúc, truy xuất đúng
bằng chứng cho nhiều kiểu câu hỏi và tạo câu trả lời có thể kiểm chứng bằng dẫn
chứng nguồn.

\section{Nhu cầu và kịch bản sử dụng}
\label{section:user-needs}

Người dùng mục tiêu của hệ thống là những người thường xuyên cần tra cứu thông
tin trong các tài liệu PDF dài hoặc có cấu trúc bán hình thức. Trong môi trường giáo
dục, người dùng có thể cần hỏi về quy chế đào tạo, điều kiện tốt nghiệp, cách tính
điểm hoặc quy định học vụ. Trong môi trường tổ chức, người dùng có thể tra cứu
quy trình nội bộ, tài liệu vận hành, hướng dẫn nghiệp vụ hoặc chính sách áp dụng
cho một trường hợp cụ thể.

Điểm chung của các kịch bản này là người dùng không chỉ cần kết quả nhanh mà
còn cần biết câu trả lời được lấy từ đâu. Vì vậy, hệ thống cần hỗ trợ hỏi đáp bằng
ngôn ngữ tự nhiên, truy xuất đúng đoạn thông tin liên quan và cung cấp dẫn chứng
nguồn. Đối với các tài liệu mang tính quy định, dẫn chứng là yêu cầu quan trọng vì
người dùng thường cần kiểm tra lại điều khoản hoặc bảng nguồn trước khi sử dụng
kết quả.

\section{Khảo sát các hướng tiếp cận liên quan}
\label{section:related-approaches}

Các hệ thống truy xuất và hỏi đáp trên tài liệu thường kết hợp nhiều hướng tiếp
cận. Truy xuất theo từ khóa như BM25 có ưu điểm nhanh, dễ triển khai và bám sát
thuật ngữ trong tài liệu \cite{robertson2009bm25}. Truy xuất ngữ nghĩa biểu diễn
câu hỏi và đoạn văn bản trong không gian vector, nhờ đó có thể tìm các đoạn liên
quan ngay cả khi câu hỏi không trùng từ khóa với tài liệu
\cite{karpukhin2020dpr,reimers2019sentencebert}. Tuy nhiên, trong tài liệu quy
định, chỉ dùng một trong hai hướng có thể chưa đủ: từ khóa có thể bỏ sót cách diễn
đạt khác, còn truy xuất ngữ nghĩa có thể lấy nhầm đoạn gần nghĩa nhưng sai điều
kiện.

Hướng hỏi đáp tăng cường bằng truy xuất sử dụng các đoạn đã truy xuất làm ngữ
cảnh cho bước sinh câu trả lời \cite{lewis2020rag}. Các hướng gần đây cũng nhấn
mạnh vai trò của việc kiểm tra lại quá trình truy xuất, sinh và phê bình câu trả lời
nhằm giảm nguy cơ tạo thông tin thiếu căn cứ \cite{asai2024selfrag}. Trong phạm
vi đồ án, các hướng tiếp cận này được sử dụng như nền tảng để thiết kế quy trình
truy xuất, kiểm tra bằng chứng và sinh câu trả lời có dẫn chứng.

\begin{table}[H]
\centering
\small
\renewcommand{\arraystretch}{1.2}
\caption{Tổng hợp các hướng tiếp cận liên quan và vai trò trong đồ án}
\label{tab:related-work-summary}
\begin{tabularx}{\textwidth}{
    |>{\raggedright\arraybackslash}p{0.22\textwidth}
    |>{\raggedright\arraybackslash}p{0.26\textwidth}
    |>{\raggedright\arraybackslash}X
    |>{\raggedright\arraybackslash}p{0.18\textwidth}|
}
\hline
\textbf{Hướng tiếp cận} & \textbf{Điểm mạnh} & \textbf{Hạn chế trong bài toán PDF QA} & \textbf{Nguồn} \\
\hline
Truy xuất theo từ khóa &
Nhanh, bám sát thuật ngữ và điều khoản trong tài liệu. &
Dễ bỏ sót khi câu hỏi diễn đạt khác văn bản gốc. &
\cite{robertson2009bm25} \\
\hline
Truy xuất ngữ nghĩa &
Tìm kiếm theo ý nghĩa của câu hỏi, phù hợp với câu hỏi tự nhiên. &
Có thể lấy nhầm đoạn gần nghĩa nhưng không chứa đúng điều kiện cần trả lời. &
\cite{karpukhin2020dpr,reimers2019sentencebert} \\
\hline
Kết hợp nhiều kết quả truy xuất &
Tận dụng nhiều tín hiệu xếp hạng khác nhau. &
Cần lựa chọn trọng số hoặc cơ chế hợp nhất phù hợp từng tập dữ liệu. &
\cite{cormack2009rrf} \\
\hline
Hỏi đáp tăng cường bằng truy xuất &
Sinh câu trả lời dựa trên ngữ cảnh được truy xuất từ tài liệu. &
Vẫn cần kiểm soát chất lượng bằng chứng và dẫn chứng nguồn. &
\cite{lewis2020rag} \\
\hline
Kiểm soát sinh câu trả lời &
Giảm nguy cơ tạo thông tin không được tài liệu hỗ trợ. &
Không bảo đảm loại bỏ hoàn toàn lỗi nếu bằng chứng đầu vào thiếu hoặc sai. &
\cite{asai2024selfrag,jeong2024adaptiverag} \\
\hline
Phân tích bố cục và bảng &
Giúp nhận diện vùng nội dung, bảng, hàng, cột và ô. &
Cần kết hợp với văn bản hoặc OCR để phục vụ truy xuất và hỏi đáp. &
\cite{pfitzmann2022doclaynet,smock2022pubtables} \\
\hline
Nhận dạng ký tự quang học &
Cho phép xử lý tài liệu được quét từ ảnh. &
Kết quả phụ thuộc chất lượng ảnh, độ nghiêng và nhiễu của tài liệu. &
\cite{du2020ppocr} \\
\hline
\end{tabularx}
\end{table}

\section{Phân tích yêu cầu của hệ thống}
\label{section:system-requirements}

Dựa trên bài toán và các kịch bản sử dụng đã phân tích, hệ thống cần đáp ứng các
yêu cầu chức năng và phi chức năng sau. Các yêu cầu được mã hóa để thuận tiện
cho việc liên kết với thiết kế ở Chương 4, Chương 5 và thực nghiệm ở Chương 6.

\subsection{Yêu cầu chức năng}

\begin{table}[H]
\centering
\small
\renewcommand{\arraystretch}{1.2}
\caption{Yêu cầu chức năng của hệ thống}
\label{tab:functional-requirements}
\begin{tabularx}{\textwidth}{
    |>{\centering\arraybackslash}p{0.12\textwidth}
    |>{\raggedright\arraybackslash}p{0.26\textwidth}
    |>{\raggedright\arraybackslash}X|
}
\hline
\textbf{Mã} & \textbf{Yêu cầu} & \textbf{Mô tả} \\
\hline
FR-01 & Tiếp nhận tài liệu PDF &
Hệ thống cho phép xử lý tài liệu PDF đầu vào và xác định thông tin cơ bản của
tài liệu. \\
\hline
FR-02 & Trích xuất nội dung &
Hệ thống trích xuất văn bản, bảng và metadata nguồn như tài liệu, trang, vùng
nội dung và loại block. \\
\hline
FR-03 & Định tuyến theo vùng nội dung &
Hệ thống có khả năng lựa chọn cách xử lý phù hợp cho vùng văn bản, bảng hoặc
vùng ảnh thay vì xử lý cả trang như một khối duy nhất. \\
\hline
FR-04 & Chia đoạn và lập chỉ mục &
Hệ thống chuyển nội dung đã trích xuất thành các đơn vị truy xuất và xây dựng
chỉ mục phục vụ tìm kiếm. \\
\hline
FR-05 & Truy xuất bằng chứng &
Hệ thống nhận câu hỏi bằng ngôn ngữ tự nhiên và trả về các đoạn thông tin liên
quan nhất. \\
\hline
FR-06 & Hỏi đáp có căn cứ &
Hệ thống tạo câu trả lời dựa trên bằng chứng đã truy xuất và tránh tự suy đoán
khi thiếu thông tin. \\
\hline
FR-07 & Dẫn chứng nguồn &
Câu trả lời cần đi kèm dẫn chứng tới tài liệu, trang, đoạn văn, bảng, hàng hoặc ô
bảng khi có thể. \\
\hline
FR-08 & Đánh giá thực nghiệm &
Hệ thống hỗ trợ chạy benchmark để đánh giá ingest, trích xuất bảng, truy xuất,
hỏi đáp và dẫn chứng. \\
\hline
\end{tabularx}
\end{table}

\subsection{Yêu cầu phi chức năng}

\begin{table}[H]
\centering
\small
\renewcommand{\arraystretch}{1.2}
\caption{Yêu cầu phi chức năng của hệ thống}
\label{tab:non-functional-requirements}
\begin{tabularx}{\textwidth}{
    |>{\centering\arraybackslash}p{0.12\textwidth}
    |>{\raggedright\arraybackslash}p{0.26\textwidth}
    |>{\raggedright\arraybackslash}X|
}
\hline
\textbf{Mã} & \textbf{Yêu cầu} & \textbf{Mô tả} \\
\hline
NFR-01 & Khả năng kiểm chứng &
Kết quả cần giữ được thông tin nguồn để người dùng kiểm tra lại trong PDF gốc. \\
\hline
NFR-02 & Tính có căn cứ &
Câu trả lời phải bám vào bằng chứng được truy xuất và hạn chế sinh thông tin
không có trong tài liệu. \\
\hline
NFR-03 & Tính mô-đun &
Các tầng ingest, chunking, retrieval và QA cần được tách riêng để dễ kiểm thử,
thay thế và đánh giá. \\
\hline
NFR-04 & Hiệu năng chấp nhận được &
Thời gian truy xuất và trả lời cần phù hợp với kịch bản hỏi đáp tương tác trên
một hoặc một nhóm tài liệu. \\
\hline
NFR-05 & Khả năng mở rộng đánh giá &
Hệ thống cần cho phép thêm benchmark hoặc cấu hình mới để phục vụ so sánh và
phân tích lỗi. \\
\hline
\end{tabularx}
\end{table}

\section{Phạm vi dữ liệu và thực nghiệm}
\label{section:data-and-experiment-scope}

Phạm vi chính của đồ án là các tài liệu PDF có lớp văn bản và cấu trúc bán hình
thức, ví dụ quy chế, quy định, thông tư, sổ tay hướng dẫn, báo cáo và tài liệu
nghiệp vụ. Đây là nhóm tài liệu thường có nhu cầu tra cứu cao và cần câu trả lời
có dẫn chứng. Đồ án có khảo sát thêm tài liệu được quét từ ảnh và tài liệu chứa
bảng để đánh giá giới hạn của pipeline, nhưng chưa đặt mục tiêu xử lý hoàn chỉnh
mọi dạng PDF.

Các trường hợp ngoài phạm vi chính gồm tài liệu ảnh chất lượng thấp, biểu đồ yêu
cầu phân tích hình ảnh chuyên sâu, công thức phức tạp, bảng cần tái dựng chính
xác tuyệt đối và câu hỏi đòi hỏi suy luận số học nhiều bước. Vì vậy, các kết quả
thực nghiệm trong đồ án chỉ được diễn giải trong phạm vi dữ liệu và cấu hình được
nêu ở Chương 6.

\section{Tiêu chí đánh giá}
\label{section:evaluation-criteria-chapter2}

Đồ án đánh giá hệ thống theo nhiều tầng thay vì chỉ đo câu trả lời cuối cùng. Ở
tầng xử lý tài liệu, tiêu chí quan trọng là khả năng trích xuất đúng nội dung, nhận
diện vùng bảng, giữ thứ tự đọc và bảo toàn metadata nguồn. Ở tầng truy xuất, hệ
thống được đánh giá theo khả năng tìm đúng bằng chứng và xếp bằng chứng phù
hợp ở vị trí cao. Ở tầng hỏi đáp, các tiêu chí chính là độ khớp câu trả lời, độ khớp
bằng chứng, tỷ lệ câu trả lời có căn cứ và tỷ lệ sinh thông tin thiếu căn cứ. Đối với
dữ liệu bảng, đồ án đánh giá thêm khả năng truy xuất đúng bảng, hàng, cột hoặc ô
liên quan.

Các định nghĩa metric và công thức được trình bày tập trung tại Chương 3; quy
ước tính điểm theo từng tập dữ liệu được nêu trong Chương 6 để tránh lặp lại giữa
phần khảo sát và phần thực nghiệm.

\section{Liên kết yêu cầu với thiết kế và thực nghiệm}
\label{section:requirements-traceability}

Bảng~\ref{tab:requirements-traceability} liên kết các yêu cầu đã xác định với phần
thiết kế ở Chương 4, Chương 5 và phần thực nghiệm ở Chương 6. Bảng này giúp đảm bảo mỗi
yêu cầu quan trọng đều có thành phần triển khai và tiêu chí đánh giá tương ứng.

\begin{table}[H]
\centering
\scriptsize
\renewcommand{\arraystretch}{1.2}
\caption{Liên kết yêu cầu với thiết kế và thực nghiệm}
\label{tab:requirements-traceability}
\begin{tabularx}{\textwidth}{
    |>{\centering\arraybackslash}p{0.10\textwidth}
    |>{\raggedright\arraybackslash}p{0.25\textwidth}
    |>{\raggedright\arraybackslash}p{0.24\textwidth}
    |>{\raggedright\arraybackslash}X|
}
\hline
\textbf{Mã} & \textbf{Yêu cầu} & \textbf{Thiết kế liên quan} & \textbf{Thực nghiệm liên quan} \\
\hline
FR-01, FR-02 & Tiếp nhận và trích xuất PDF &
Tầng ingest tài liệu PDF &
Mục~\ref{subsection:ingest-results} và Mục~\ref{subsection:region-routing-results} \\
\hline
FR-03 & Định tuyến theo vùng nội dung &
Region-level routing, metadata vùng &
Mục~\ref{subsection:region-routing-results} \\
\hline
FR-04 & Chia đoạn và lập chỉ mục &
Chunking văn bản, table-aware chunking, indexing &
Mục~\ref{subsection:table-aware-results} và Mục~\ref{subsection:retrieval-results} \\
\hline
FR-05 & Truy xuất bằng chứng &
BM25, dense retrieval, hybrid retrieval, reranking &
Mục~\ref{subsection:retrieval-results} \\
\hline
FR-06, FR-07 & Hỏi đáp có căn cứ và dẫn chứng &
Evidence checker, answer generator, citation &
Mục~\ref{section:end-to-end-results} và Mục~\ref{section:ablation-results} \\
\hline
FR-08, NFR-05 & Đánh giá thực nghiệm &
Script benchmark, cấu hình ablation &
Chương 6, các mục thực nghiệm theo tầng và ablation \\
\hline
NFR-01, NFR-02 & Khả năng kiểm chứng và tính có căn cứ &
Metadata nguồn, citation target, cell-level evidence &
Mục~\ref{subsection:table-aware-results}, Mục~\ref{section:end-to-end-results} \\
\hline
NFR-03, NFR-04 & Tính mô-đun và hiệu năng &
Tổ chức pipeline theo tầng, lưu index tách biệt &
Mục~\ref{subsection:region-routing-results} và Mục~\ref{section:error-analysis} \\
\hline
\end{tabularx}
\end{table}

\section{Kết chương}

Chương này đã xác định rõ bài toán truy xuất và hỏi đáp trên tài liệu PDF, các
thách thức chính của dữ liệu PDF, nhu cầu sử dụng, các hướng tiếp cận liên quan
và phạm vi thực nghiệm của đồ án. Điểm khác biệt so với phần giới thiệu là các
yêu cầu đã được mã hóa thành FR/NFR và liên kết với thiết kế, thực nghiệm ở các
chương sau. Đây là cơ sở để Chương 3 trình bày các nền tảng lý thuyết và công
nghệ được sử dụng, trước khi Chương 4 đi vào thiết kế hệ thống cụ thể.

\end{document}
